\section{Coherent ideal powers and the original Rees degrees}

\begin{lemma}[The ideal-power Čech complex]
\label{lem:fc-rees-cech-ideal-power-sections}
\leanok
\lean{EnvelopeProperty.FcReesCechIdealPowerSections.baseScalar, EnvelopeProperty.FcReesCechIdealPowerSections.chartLinearIso, EnvelopeProperty.FcReesCechIdealPowerSections.chartLinearIso_naturality, EnvelopeProperty.FcReesCechIdealPowerSections.factorIso, EnvelopeProperty.FcReesCechIdealPowerSections.termIso, EnvelopeProperty.FcReesCechIdealPowerSections.termIso_projection, EnvelopeProperty.FcReesCechIdealPowerSections.factorIso_face, EnvelopeProperty.FcReesCechIdealPowerSections.termIso_coface, EnvelopeProperty.FcReesCechIdealPowerSections.termIso_differential, EnvelopeProperty.FcReesCechIdealPowerSections.complexIso, EnvelopeProperty.FcReesCechIdealPowerSections.homologyIso}
\uses{lem:fc-rees-cech-degree-complex}
Let $X\to\operatorname{Spec}R$ be separated, let $I\subseteq R$ be an
ideal, and let $(U_i)$ be a finite linearly ordered family of affine opens
of $X$.  For every $d\geq0$, the ordered Čech complex of the ideal
sheaf $I^d\mathcal O_X$, regarded as an $R$-linear complex through the
original structural map, is canonically isomorphic to the degree-$d$
complex obtained from the schematic Rees family.  Its map on every
intersection is the identification with
$(I\mathcal O_X(U_{i_0}\cap\cdots\cap U_{i_q}))^d$; in particular the
comparison preserves every restriction and every coface.  It therefore
identifies their cohomology in every degree.
\end{lemma}
\begin{proof}
On an affine open $V$, the sheaf associated to the extended ideal power
has sections equal to that ideal power inside $\mathcal O_X(V)$.  The
structural map from $R$ to these sections is the original coefficient
map, so this identification is $R$-linear.  If $V\subseteq W$ are affine,
both maps from sections over $W$ to the ideal power over $V$ are the
restriction in the structure sheaf, hence agree on every section.
Every intersection of members of the affine family is affine because
$X$ is separated.  Taking the product of the affine identifications over
all increasing intersection indices gives an isomorphism in each Čech
degree.  A coface is an intersection projection followed by precisely
one of these restrictions; the restriction equality therefore proves
that the whole coface square commutes.  Taking the alternating finite
sum proves compatibility with the differential.  These termwise
isomorphisms give an isomorphism of complexes.  Restriction of scalars
along the affine section isomorphism is exact, so passage to cohomology
gives the asserted original $R$-linear cohomology comparison.
\end{proof}
